\documentclass[12pt,a4paper]{article}
\usepackage[utf8]{inputenc}
\usepackage{amsmath,amssymb,amsthm}
\usepackage{mathtools}
\usepackage{geometry}
\geometry{margin=1in}

\newtheorem{theorem}{Theorem}[section]
\newtheorem{lemma}[theorem]{Lemma}
\newtheorem{definition}[theorem]{Definition}
\newtheorem{remark}[theorem]{Remark}
\newtheorem{conjecture}{Conjecture}[section]

\title{Reformulating the Riemann Hypothesis via a Greedy Harmonic Condition:
A Research Roadmap}
\author{}
\date{}

\begin{document}

\maketitle

\begin{abstract}
We present a reformulation of the Riemann Hypothesis (RH) in terms of a greedy condition on the partial sums of the Dirichlet eta function. After establishing the equivalence, we outline a concrete research program toward a possible proof. This program consists of four major steps: (1) proving that every zero satisfies the greedy condition if and only if it lies on the critical line, (2) developing effective bounds for the argument of partial sums, (3) using modern exponential‑sum techniques to control off‑line zeros, and (4) applying the greedy condition to derive a contradiction for any zero with real part different from $1/2$. We also suggest computational experiments and connections to existing literature. The goal is not to claim a proof, but to provide a clear, testable pathway for future work.
\end{abstract}

\section{Introduction and Recap}

In the companion paper we proved that the Riemann Hypothesis is equivalent to the statement that every non‑trivial zero $s=\sigma+it$ of the Dirichlet eta function
\[
\eta(s)=\sum_{n=1}^\infty (-1)^{n-1}n^{-s}
\]
satisfies the \textbf{greedy condition}:
\[
\operatorname{Re}\bigl(\overline{S_{N-1}(s)}\,w_N(s)\bigr)\le 0\qquad\forall N\ge1,
\]
where $S_{N-1}(s)=\sum_{n=1}^{N-1}(-1)^{n-1}n^{-s}$ and $w_N(s)=N^{-s}$. The present article outlines a research program aimed at proving that this greedy condition forces $\sigma=1/2$. If successful, this would establish RH. We emphasise that no proof is given here; instead we describe concrete mathematical steps that, if completed, would yield a proof.

\section{Step 1: Strengthening the Greedy Condition to an Energy Inequality}

The greedy condition is equivalent to
\[
|S_N(s)|^2 \le |S_{N-1}(s)|^2 + |w_N(s)|^2 - 2|S_{N-1}(s)||w_N(s)|\sin^2(\phi_N/2),
\]
where $\phi_N$ is the angle between $S_{N-1}$ and $w_N$. A necessary condition for the inequality to hold for all $N$ is that the angle $\phi_N$ lies in $[-\pi/2,\pi/2]$ modulo $2\pi$. We propose to study the cumulative energy
\[
E(N)=|S_N(s)|^2 + \sum_{k=N+1}^\infty |w_k(s)|^2.
\]
If the greedy condition holds and $\eta(s)=0$, one can show that $E(N)$ is non‑increasing. A future proof would require showing that for $\sigma\neq1/2$, $E(N)$ cannot be bounded below by zero unless a contradiction arises (e.g., $E(N)$ would become negative).

\section{Step 2: Asymptotic Control of the Argument of Partial Sums}

The key analytic difficulty is to understand the argument of $S_{N-1}(s)$ for large $N$ when $\eta(s)=0$. The Riemann–Siegel formula provides an asymptotic expansion:
\[
S_N(s) = \frac{1}{\sqrt{2\pi}} N^{-1/4} e^{i(\theta(t)-\sqrt{2\pi N})} + O(N^{-3/4}),\quad s=\frac12+it,
\]
but for $\sigma\neq1/2$ the formula changes. A research task is to derive a uniform asymptotic expansion for $S_N(s)$ for all $0<\sigma<1$, valid for $N$ in the range $N\approx \sqrt{t}$. Using the approximate functional equation (Titchmarsh, Theorem 4.15), one can write
\[
S_N(s) = \eta(s) + \chi(s)\sum_{n\le N} \frac{(-1)^{n-1}}{n^{1-s}} + \text{error},
\]
where $\chi(s)=2^{s}\pi^{s-1}\sin(\pi s/2)\Gamma(1-s)$. For $\eta(s)=0$, the first term vanishes. The second term can be estimated by stationary phase. The leading term’s magnitude is
\[
\left|\chi(s)\sum_{n\le N} \frac{(-1)^{n-1}}{n^{1-s}}\right| \asymp N^{\sigma-1/2}\cdot |\chi(s)|.
\]
For $\sigma>1/2$ this decays; for $\sigma<1/2$ it grows. The greedy condition imposes a sign constraint on the real part of $\overline{S_{N-1}}w_N$. A rigorous analysis would show that this constraint can be satisfied for all $N$ only if the leading term is purely imaginary (so that its real part with $w_N$ is zero up to error). This happens when the argument of the leading term equals $\pi/2$ modulo $\pi$, which is equivalent to a condition on $t$ that defines the zeros. Moreover, the magnitude must not grow, forcing $\sigma=1/2$. This is the central analytic challenge.

\section{Step 3: Modern Exponential‑Sum Methods}

Classical stationary phase gives asymptotic formulas with error terms that are not uniform in $t$ for all ranges. To obtain effective bounds (with explicit constants) that work for all $N$, one can use modern decoupling theory (Bourgain, Demeter, Guth) and the recent work of Matomäki, Radziwiłł, and Tao on short‑interval discorrelation. Specifically, one needs to estimate sums of the form
\[
\sum_{n=1}^N (-1)^{n-1} n^{-\sigma} e^{-it\ln n}
\]
for large $N$ and $t$ that are not necessarily aligned with the Riemann–Siegel transition point. The greedy condition is a pointwise inequality for every $N$, not just asymptotically. To prove that it forces $\sigma=1/2$, one could use a variant of the “energy method” combined with the fact that the set of $t$ for which the inequality holds for all $N$ must have measure zero unless $\sigma=1/2$ (by density of the phases). A promising approach is to use the following lemma:

\begin{lemma}[Density of phases]
If $t/(2\pi)$ is irrational, the sequence $\{t\ln n \mod 2\pi\}_{n\ge1}$ is dense in $[0,2\pi)$.
\end{lemma}

Then, for a fixed $\sigma\neq1/2$, one can show that the greedy condition would imply that a certain Fourier series vanishes, forcing $t$ to be such that the associated Dirichlet series has a zero, but also forcing a contradiction via the functional equation. This line of reasoning is reminiscent of the uniqueness of the critical line in the theory of Dirichlet series with multiplicative coefficients.

\section{Step 4: A Contradiction for Off‑Line Zeros}

Assume that there exists a zero $s=\sigma+it$ with $\sigma\neq1/2$ and $0<\sigma<1$. Then by the functional equation, $1-s$ is also a zero. One can then consider the product $\eta(s)\eta(1-s)$. Using the greedy condition on both zeros, one might derive a contradiction using the fact that the greedy condition implies a certain inequality on the real parts of the partial sums that is incompatible with the symmetry $s\leftrightarrow 1-s$. More concretely, the greedy condition for $s$ gives a constraint on the argument of $S_N(s)$; the same for $1-s$ gives a constraint on the argument of $S_N(1-s)$. The two constraints together would force $\sigma=1/2$ by an elementary argument on the angle between $S_N(s)$ and $S_N(1-s)$. This idea is yet to be developed.

\section{Computational Verification and Heuristics}

As a preliminary step, we encourage extensive numerical testing of the greedy condition for known zeros (on the critical line) and for randomly chosen $s$ with $\sigma\neq1/2$. The condition can be checked efficiently because the series converges conditionally; one can compute partial sums up to $N\approx 10^6$ and verify the inequality. For off‑line points (e.g., $s=0.6+14.1347i$), we expect the inequality to fail for some $N$. Such numerical experiments would build confidence that the greedy condition indeed characterises the critical line. Moreover, they could reveal a minimal counterexample $N$ that grows with $t$, giving hints for analytic estimates.

\section{Conclusion: A Research Roadmap}

We have outlined a four‑step research program:

\begin{enumerate}
    \item Derive an energy inequality from the greedy condition and analyse its monotonicity.
    \item Obtain uniform asymptotic expansions for the partial sums of $\eta(s)$ for all $\sigma\in(0,1)$, using the approximate functional equation and modern stationary phase methods.
    \item Apply decoupling and discorrelation estimates to bound the error terms and to show that the greedy condition forces the leading term to be purely imaginary and the magnitude to decay, implying $\sigma=1/2$.
    \item Use the functional equation to derive a contradiction for any hypothetical zero off the critical line, thereby proving RH.
\end{enumerate}

This roadmap is ambitious but each step is plausible with current analytic number theory tools. The reformulation via the greedy condition does not simplify the problem, but it provides a new geometric viewpoint that may be amenable to combinatorial or Fourier‑analytic attacks. We hope that this research program inspires further work.

\begin{thebibliography}{9}
\bibitem{titchmarsh} E. C. Titchmarsh, \textit{The Theory of the Riemann Zeta Function}, Oxford University Press, 1986.
\bibitem{edwards} H. M. Edwards, \textit{Riemann's Zeta Function}, Dover, 2001.
\bibitem{bourgain} J. Bourgain, C. Demeter, L. Guth, ``Proof of the main conjecture in Vinogradov's mean value theorem for degrees higher than three'', Ann. of Math. 184 (2016), 633–682.
\bibitem{mrt} K. Matomäki, M. Radziwiłł, T. Tao, ``An averaged form of Chowla's conjecture'', Algebra Number Theory 15 (2021), 2167–2240.
\end{thebibliography}

\end{document}